% Packages, typography, macros, and theorem environments for Relational Superposition Theory

\usepackage[utf8]{inputenc}
\usepackage[T1]{fontenc}
\usepackage{lmodern}
\usepackage{microtype}
\usepackage[top=1in,bottom=1in,left=1in,right=1in,headheight=23pt]{geometry}

\usepackage{amsmath,amssymb,amsthm,mathtools}
\usepackage{booktabs,tabularx,multirow,array}
\usepackage{enumitem}
\usepackage{setspace}
\usepackage{xcolor}
\usepackage{graphicx}
\usepackage{float}
\usepackage{placeins}
\usepackage{needspace}
\usepackage{caption}
\usepackage{subcaption}
\usepackage{listings}
\usepackage{fancyhdr}
\usepackage{titlesec}
\usepackage{tocloft}
\usepackage[numbers,sort&compress]{natbib}
\usepackage{xurl}
\usepackage{tikz}
\usepackage{pgfplots}
\pgfplotsset{compat=1.18}
\usetikzlibrary{arrows.meta,positioning,shapes.geometric,calc,fit,backgrounds}

% Color palette: restrained editorial ink and slate
\definecolor{inkblack}{HTML}{141416}
\definecolor{slateink}{HTML}{334155}
\definecolor{mutedgray}{HTML}{64748B}
\definecolor{rulegray}{HTML}{CBD5E1}
\definecolor{accentink}{HTML}{2B3A67}
\definecolor{plumink}{HTML}{5B3E6B}
\definecolor{softtint}{HTML}{F8FAFC}
\definecolor{codebg}{HTML}{F6F7F9}

\setstretch{1.10}
\widowpenalty=10000
\clubpenalty=10000
\setlength{\parskip}{0.35em}
\setlength{\parindent}{1.2em}

% Running headers and footers
\pagestyle{fancy}
\fancyhf{}
\fancyhead[L]{\small\color{mutedgray}\nouppercase{\leftmark}}
\fancyhead[R]{\small\color{mutedgray}\thepage}
\renewcommand{\headrulewidth}{0.4pt}
\renewcommand{\headrule}{\hbox to\headwidth{\color{rulegray}\leaders\hrule height \headrulewidth\hfill}}
\renewcommand{\footrulewidth}{0pt}

% Section formatting
\titleformat{\section}
  {\needspace{5\baselineskip}\Large\bfseries\color{inkblack}}
  {\thesection}{0.8em}{}[\vspace{0.25em}{\color{rulegray}\titlerule[0.5pt]}]
\titleformat{\subsection}
  {\needspace{4\baselineskip}\large\bfseries\color{inkblack}}
  {\thesubsection}{0.7em}{}
\titleformat{\subsubsection}
  {\needspace{3\baselineskip}\normalsize\bfseries\color{slateink}}
  {\thesubsubsection}{0.6em}{}

% Caption styling
\captionsetup{
  font=small,
  labelfont={bf,color=inkblack},
  textfont={color=slateink},
  margin=1em
}

% Listings styling for Lean 4 and pseudocode
\lstdefinelanguage{Lean4}{
  morekeywords={import,namespace,end,structure,inductive,def,noncomputable,theorem,lemma,proposition,example,where,let,in,if,then,else,match,with,by,have,show,calc,fun,forall,exists,open,section,variable,private,protected,deriving,instance,class},
  sensitive=true,
  morecomment=[l]{--},
  morecomment=[s]{/-}{-/},
  morestring=[b]"
}

\lstset{
  basicstyle=\ttfamily\footnotesize,
  backgroundcolor=\color{codebg},
  frame=single,
  rulecolor=\color{rulegray},
  breaklines=true,
  columns=fullflexible,
  keepspaces=true,
  showstringspaces=false,
  xleftmargin=0.8em,
  xrightmargin=0.8em,
  aboveskip=0.8em,
  belowskip=0.8em
}

% Hyperref and cleveref (cleveref must follow hyperref and precede theorem definitions)
\usepackage[
  colorlinks=true,
  linkcolor=accentink,
  citecolor=accentink,
  urlcolor=accentink,
  bookmarksnumbered=true,
  pdftitle={The Rank-Role Uncertainty Principle: A Theory of Relational Superposition},
  pdfauthor={Ratnadwip Sarkar},
  pdfsubject={Relational Superposition Theory},
  pdfkeywords={relational superposition, situationship, rank, role, uncertainty, collapse, decoherence, optionality}
]{hyperref}
\usepackage[nameinlink,noabbrev]{cleveref}

% Theorem environments (numbered within section for clean cross-referencing)
\theoremstyle{plain}
\newtheorem{theorem}{Theorem}[section]
\newtheorem{principle}[theorem]{Principle}
\newtheorem{proposition}[theorem]{Proposition}
\newtheorem{lemma}[theorem]{Lemma}
\newtheorem{corollary}[theorem]{Corollary}
\newtheorem{conjecture}[theorem]{Conjecture}

\theoremstyle{definition}
\newtheorem{definition}[theorem]{Definition}
\newtheorem{assumption}[theorem]{Assumption}
\newtheorem{example}[theorem]{Example}

\theoremstyle{remark}
\newtheorem{remark}[theorem]{Remark}

% Cleveref names
\crefname{theorem}{Theorem}{Theorems}
\crefname{principle}{Principle}{Principles}
\crefname{proposition}{Proposition}{Propositions}
\crefname{lemma}{Lemma}{Lemmas}
\crefname{corollary}{Corollary}{Corollaries}
\crefname{conjecture}{Conjecture}{Conjectures}
\crefname{definition}{Definition}{Definitions}
\crefname{assumption}{Assumption}{Assumptions}
\crefname{example}{Example}{Examples}
\crefname{remark}{Remark}{Remarks}

% Core mathematical macros
\newcommand{\RST}{\textsc{rst}}
\newcommand{\ket}[1]{\lvert #1 \rangle}
\newcommand{\bra}[1]{\langle #1 \rvert}
\newcommand{\braket}[2]{\langle #1 \mid #2 \rangle}
\newcommand{\ketbra}[2]{\lvert #1 \rangle\langle #2 \rvert}
\newcommand{\expect}[2]{\langle #1 \rangle_{#2}}
\newcommand{\var}[2]{\operatorname{Var}_{#2}(#1)}
\newcommand{\std}[2]{\Delta_{#2} #1}
\newcommand{\comm}[2]{[#1, #2]}
\newcommand{\opR}{\hat{R}}
\newcommand{\opL}{\hat{L}}
\newcommand{\opI}{\hat{I}}
\newcommand{\poke}{\mathcal{P}}
\newcommand{\kappaS}{\kappa}
\newcommand{\roleset}{\mathcal{R}}
\newcommand{\partnerset}{\mathcal{G}}
\newcommand{\vspaceR}{\mathcal{V}_{\mathcal{R}}}
\newcommand{\vspaceG}{\mathcal{V}_{\mathcal{G}}}
\newcommand{\vspaceTot}{\mathcal{V}}
\newcommand{\Cnum}{\mathbb{C}}
\newcommand{\Rnum}{\mathbb{R}}
\newcommand{\Qnum}{\mathbb{Q}}
\newcommand{\Znum}{\mathbb{Z}}
\newcommand{\Nnum}{\mathbb{N}}
\newcommand{\Prob}{\mathbb{P}}
\newcommand{\E}{\mathbb{E}}
\newcommand{\norm}[1]{\lVert #1 \rVert}
\newcommand{\abs}[1]{\lvert #1 \rvert}
\newcommand{\diag}{\operatorname{diag}}
\newcommand{\tr}{\operatorname{tr}}
\newcommand{\supp}{\operatorname{supp}}
\newcommand{\RePart}{\operatorname{Re}}
\newcommand{\ImPart}{\operatorname{Im}}
